\section{The Python Code}

We wrote two Python functions during the project. The first one, \texttt{J\_0},
plays the game step by step. The second one, \texttt{J\_1}, uses the formula of
Section~4. In this section we explain how they work, we compare their cost,
and we give a faster version of the formula together with the program we used
to check all our results.

\Needspace{23\baselineskip}
\subsection{\texorpdfstring{\texttt{J\_0}: a simulation of the game}{J\_0: a simulation of the game}}

\begin{lstlisting}[language=iPython, caption={The function \texttt{J\_0}, which simulates the game.}]
def J_0(n):
    if n <= 0:
        return "Please enter a positive integer for n"
    else:
        people = [i for i in range(1, n + 1)]
        print(people)
        while len(people) > 1:
            if len(people) % 2 == 0:  # Even case
                for i in range(len(people)//2):
                    h = people.pop(i + 1)
                    print(people, "\t", h, "is eliminated")
            else:  # Odd case
                for i in range(len(people)//2):
                    h = people.pop(i + 1)
                    print(people, "\t", h, "is eliminated")
                people.insert(0, people.pop(len(people) - 1))
                print(people)
        return f"So {people[0]} is the only survivor"
\end{lstlisting}

\textbf{The list.} The circle is a Python list, \texttt{people}, which starts
as $[1,2,\dots,n]$. The person at index $0$ is the next one to be skipped. If
$n\leq0$ the function returns a message, since the game makes no sense.

\textbf{One pass over the list.} Each turn of the \texttt{while} loop goes once
through the list and eliminates every second person. The instruction
\texttt{people.pop(i + 1)} removes the person at index $i+1$. This index is
the right one because the list shrinks while we go: after $i$ eliminations,
$i$ people have left the front part of the list, so the person who was at
index $2i+1$ is now at index $i+1$. For $n=6$:
\[
[1,2,3,4,5,6]\xrightarrow{\ \text{pop}(1)\ }[1,3,4,5,6]
\xrightarrow{\ \text{pop}(2)\ }[1,3,5,6]
\xrightarrow{\ \text{pop}(3)\ }[1,3,5].
\]
With $2k$ people the pass makes $k$ eliminations and ends exactly at the end
of the list, so the next pass starts again with the person at index $0$, who
is skipped.

\textbf{The odd case.} With $2k+1$ people the pass also makes $k$
eliminations, but the last person of the list is skipped at the end of the
pass. The next one to be eliminated is therefore the first person of the list.
The instruction
\begin{center}
\texttt{people.insert(0, people.pop(len(people) - 1))}
\end{center}
moves the last person to the front. Now he is at index $0$ and is skipped, and
the person after him is eliminated, as in the game. For instance
$[1,5,9]$ becomes $[9,1,5]$, and the next pass eliminates $1$.

\textbf{Output for $n=10$.} The function prints the list after each
elimination:
\Needspace{18\baselineskip}
\begin{lstlisting}[numbers=none, caption={Output of \texttt{J\_0(10)}.}]
[1, 2, 3, 4, 5, 6, 7, 8, 9, 10]
[1, 3, 4, 5, 6, 7, 8, 9, 10]   2 is eliminated
[1, 3, 5, 6, 7, 8, 9, 10]      4 is eliminated
[1, 3, 5, 7, 8, 9, 10]         6 is eliminated
[1, 3, 5, 7, 9, 10]            8 is eliminated
[1, 3, 5, 7, 9]                10 is eliminated
[1, 5, 7, 9]                   3 is eliminated
[1, 5, 9]                      7 is eliminated
[9, 1, 5]
[9, 5]                         1 is eliminated
[5, 9]
[5]                            9 is eliminated
So 5 is the only survivor
\end{lstlisting}
The people are eliminated in the order $2,4,6,8,10,3,7,1,9$, which is the order
of the statement. The passes of the program are not exactly our rounds: in a
circle with an odd number of people, we counted the elimination of the first
person in the current round (Section~\ref{sec:conventions}), while the program
does it at the start of the next pass. This changes nothing to the order of
the eliminations or to the survivor.

\Needspace{19\baselineskip}
\subsection{\texorpdfstring{\texttt{J\_1}: the formula with a slow search}{J\_1: the formula with a slow search}}

\begin{lstlisting}[language=iPython, caption={The function \texttt{J\_1}, which uses the formula $J(n)=2\ell+1$.}]
def J_1(n):
    m = 0
    if n == 1:
        return 1
    elif n <= 0:
        return "Please enter a positive integer for n"
    else:
        # Find the greatest m and the corresponding
        # remainder l such that n = 2^m + l
        for j in reversed(range(n)):
            for i in range(n):
                if n == pow(2, i) + j:
                    m, l = i, j
        return 1 + 2*l
\end{lstlisting}

\textbf{Finding $m$ and $\ell$.} The function tries every pair $(i,j)$ with
$0\leq i,j<n$ and keeps the pairs with $n=2^i+j$. The outer loop takes $j$ in
decreasing order, so the last pair kept has the smallest possible $j$. A
smaller remainder means a larger power of $2$, so the last pair kept is the one
where $2^i$ is the largest power of $2$ not bigger than $n$. This is the
decomposition of Observation~\ref{obs:decomp}, and the function returns
$2\ell+1$.

\textbf{Example.} For $n=10$, the pairs found are $(0,9)$, $(1,8)$, $(2,6)$
and $(3,2)$, since $10=1+9=2+8=4+6=8+2$. The last one is $(3,2)$, so
$m=3$, $\ell=2$, and \texttt{J\_1(10)} returns $5$, like \texttt{J\_0(10)}.

\subsection{The cost of \texttt{J\_0} and \texttt{J\_1}}

\textbf{\texttt{J\_0}.} The game has $n-1$ eliminations. Each \texttt{pop}
may move all the people after the removed one, so one elimination can cost up
to $n$ operations, and the whole game about $n^2$. On top of that, the
function prints the whole list after each elimination, which is also of order
$n^2$ characters. The memory used is of order $n$, for the list.

\textbf{\texttt{J\_1}.} The two loops test $n\times n=n^2$ pairs. Moreover,
\texttt{pow(2, i)} is a number with $i+1$ digits in base $2$, so for large $i$
it is not computed in one step. In practice \texttt{J\_1} is even slower than
\texttt{J\_0}: in our tests, for $n=2000$, \texttt{J\_1} took about $3$
seconds, while \texttt{J\_0} (with the printing switched off) took about a
tenth of a second.

So the formula is fast, but this way of using it is not: almost all the work
goes into the search for $m$ and $\ell$. The next subsection fixes this.

\subsection{Faster versions of the formula}

To find the largest power of $2$ not bigger than $n$, we do not need to try
every pair. We start from $p=1$ and double $p$ as long as $2p\leq n$. This
needs only about $\log_2 n$ steps: $30$ steps for $n=10^9$.

\Needspace{29\baselineskip}
\begin{lstlisting}[language=iPython, caption={Faster functions for $J(n)$ and for a game that starts with person $s$.}]
def J_2(n):
    # I look for the largest power of 2 that is not bigger than n,
    # by doubling p as long as I stay below n
    p = 1
    while 2 * p <= n:
        p = 2 * p
    l = n - p
    return 2 * l + 1


def J_binary(n):
    b = bin(n)[2:]               # for example 10 gives '1010'
    return int(b[1:] + b[0], 2)  # I move the first digit to the end


def J_start(n, s):
    # survivor when the game starts with person s (Section 4.2)
    p = 1
    while 2 * p <= n:
        p = 2 * p
    l = n - p
    if s + 2 * l <= n:
        return s + 2 * l
    return s + 2 * l - n
\end{lstlisting}

The function \texttt{J\_2} is the formula $J(n)=2\ell+1$ with this quick
search for $2^m$. The function \texttt{J\_binary} is the rule of
Theorem~\ref{thm:binary}: \texttt{bin(n)} writes $n$ in base $2$ (with the
prefix \texttt{0b}, which we remove), we move the first digit to the end, and
\texttt{int(..., 2)} converts back to base $10$. The function \texttt{J\_start}
is the formula of Theorem~\ref{thm:general}. These three functions give the
answer for $n=10^9$, and even for $n=10^{100}$, at once.

\subsection{Checking all the results}

To check the formulas, we wrote a second simulation, \texttt{play}, which
prints nothing and accepts any starting person. It uses a \texttt{deque}, a
list from which Python can remove the first element quickly.

\Needspace{24\baselineskip}
\begin{lstlisting}[language=iPython, caption={The simulation used to check every formula.}]
from collections import deque


def play(n, s=1):
    # I play the game without printing anything, to check the formulas
    circle = deque(range(1, n + 1))
    circle.rotate(-(s - 1))      # person s is now the first of the list
    while len(circle) > 1:
        circle.rotate(-1)        # the first person is skipped and goes to the back
        circle.popleft()         # the next one is eliminated
    return circle[0]


def check(N):
    for n in range(1, N + 1):
        assert play(n) == J_2(n) == J_binary(n)
        for s in range(1, n + 1):
            assert play(n, s) == J_start(n, s)
    return "all the checks passed for n = 1, ..., " + str(N)
\end{lstlisting}

In \texttt{play}, the first person of the deque is always the next one to be
skipped. \texttt{rotate(-1)} sends him to the back of the circle, and
\texttt{popleft()} eliminates the person who is now first. The function
\texttt{check(200)} compares, for every $n\leq200$, the game with
\texttt{J\_2} and \texttt{J\_binary}, and for every starting person $s$, the
game with \texttt{J\_start}. All the checks passed. We also checked that
\texttt{J\_0} and \texttt{J\_1} give the same survivor as the formula for
every $n\leq300$.

\subsection{Comparison of the four functions}

\begin{table}[!ht]
\centering
\renewcommand{\arraystretch}{1.2}
\begin{tabular}{|l|l|l|l|}
\hline
Function & Method & Number of steps & $n=10^9$\\ \hline
\texttt{J\_0} & plays the game & about $n^2$ & out of reach\\
\texttt{J\_1} & formula, search over all pairs & more than $n^2$ & out of reach\\
\texttt{J\_2} & formula, doubling search & about $\log_2 n$ & at once\\
\texttt{J\_binary} & first digit moved to the end & about $\log_2 n$ & at once\\ \hline
\end{tabular}
\caption{The four functions for $J(n)$.}
\label{tab:code}
\end{table}

\texttt{J\_0} is the one that helped us most at the start: it shows the
circle after each elimination, so we could compare it with our drawings.
\texttt{J\_1} was our first use of the formula. \texttt{J\_2} and
\texttt{J\_binary} show what a formula brings once it is used well: the answer
no longer depends on playing the game.
